\chapter{مبانی نظری و مرور کارهای پیشین}\label{ch:background}

\section{مبانی نظری}
نظریه را با رابطه‌های شماره‌دار بیاورید و با برچسب به آن‌ها ارجاع دهید:
\begin{equation}\label{eq:energy}
  E = mc^{2}.
\end{equation}
\Cref{eq:energy} به‌ترتیب فصل شماره‌گذاری می‌شود. برای نمادگذاری و یکاها از
\lr{\texttt{\textbackslash vect}}، \lr{\texttt{\textbackslash mat}} و \lr{\texttt{\textbackslash qty}}
استفاده کنید، مثلاً $\vect{F}=m\vect{a}$ و \qty{9.81}{\metre\per\second\squared}.

\begin{definitionbox}[تعریف]
تعریف جعبه‌ای برای مفهوم اصلی فصل.
\end{definitionbox}

\begin{theorem}[فیثاغورس]\label{thm:pyth}
در مثلث قائم‌الزاویه با ضلع‌های $a,b$ و وتر $c$ داریم $a^{2}+b^{2}=c^{2}$.
\end{theorem}
\begin{proof}
استاندارد است؛ \cite{knuth1984} را ببینید.
\end{proof}

\section{کارهای مرتبط}
روش‌های موجود را در یک جدول مقایسه کنید.
\begin{center}\atlrows
  \captionof{table}{مقایسه‌ی روش‌های موجود.}\label{tab:related}
  \begin{tabular}{@{}llll@{}}
    \rowcolor{atl-primary}\atlhead{روش} & \atlhead{مزیت} & \atlhead{عیب} & \atlhead{مرجع}\\
    روش الف & سریع & کم‌دقت & \cite{knuth1984} \\
    روش ب & دقیق & کند & \cite{lamport1994} \\
  \end{tabular}
\end{center}

\section{جمع‌بندی فصل}
بنویسید خواننده از این فصل چه چیزی باید به یاد داشته باشد.
